%! Polynomial Functions and Rates of Change — Unit 2, Lesson 2.2 — Group Activity
\documentclass[10pt]{article}
\usepackage{precalculus-article}
\usepackage{precalculus-boxes}
\newcommand{\TallMath}[1]{$\displaystyle #1\rule[-1.4em]{0pt}{3.2em}$}

\begin{document}

\pageheader{Unit 2, Lesson 2.2}{Group Activity}
\namepartnerperiod

\vspace{0.10in}

\begin{scenariobox}[Roller Coaster Heights]{plumlight}
A roller coaster's height (in feet) above the ground during a 60-second ride is modeled
by a polynomial function $h(t)$. A table of sample values is given below.

\vspace{6pt}
\renewcommand{\arraystretch}{1.4}
\begin{tabularx}{\linewidth}{@{} >{\bfseries}c *{9}{X} @{}}
\rowcolor{sky}
$t$ (sec) & 0 & 10 & 15 & 20 & 30 & 35 & 40 & 50 & 60 \\
\midrule
$h(t)$ (ft) & 0 & 80 & 120 & 80 & 30 & 50 & 30 & 10 & 0 \\
\end{tabularx}
\end{scenariobox}

\vspace{0.10in}

\begin{tcolorbox}[colback=white, colframe=black!40,
    title=\textbf{Tier R --- Remediate: Reading Key Features from a Table},
    fonttitle=\bfseries, arc=2mm, left=3mm, right=3mm, top=2mm, bottom=2mm]

Use the table above (no graphing needed).

\begin{enumerate}[leftmargin=1.5em, itemsep=8pt]
    \item Identify the time(s) at which the roller coaster reaches a \textbf{local maximum}
    (a high point between two lower values). List the time and height.

    \vspace{4pt}
    \blank{10cm}

    \item Identify the time(s) at which the roller coaster reaches a \textbf{local minimum}
    (a low point between two higher values). List the time and height.

    \vspace{4pt}
    \blank{10cm}

    \item What is the \textbf{global maximum} height over the entire ride?
    At what time does it occur?

    \vspace{4pt}
    \blank{10cm}

    \item The coaster returns to height 0 at $t = 0$ and $t = 60$.
    Based on the table, does the function appear to have any other zeros? If so, where?

    \vspace{4pt}
    \blank{10cm}
\end{enumerate}
\end{tcolorbox}

\vspace{0.10in}

\begin{tcolorbox}[colback=white, colframe=black!40,
    title=\textbf{Tier A --- On Level: Rates of Change and Concavity},
    fonttitle=\bfseries, arc=2mm, left=3mm, right=3mm, top=2mm, bottom=2mm]

\begin{enumerate}[leftmargin=1.5em, itemsep=8pt]
    \item Compute the AROC of $h$ over each interval. Show your work.

    \renewcommand{\arraystretch}{1.6}
    \begin{tabularx}{\linewidth}{@{} l X X @{}}
    \rowcolor{sky}
    Interval & AROC formula & AROC value \\
    \midrule
    $[0, 10]$   & $\dfrac{h(10)-h(0)}{10-0}$   & \blank{2.5cm} \\
    $[10, 15]$  & $\dfrac{h(15)-h(10)}{15-10}$  & \blank{2.5cm} \\
    $[15, 20]$  & $\dfrac{h(20)-h(15)}{20-15}$  & \blank{2.5cm} \\
    $[20, 30]$  & $\dfrac{h(30)-h(20)}{30-20}$  & \blank{2.5cm} \\
    $[30, 35]$  & $\dfrac{h(35)-h(30)}{35-30}$  & \blank{2.5cm} \\
    $[35, 40]$  & $\dfrac{h(40)-h(35)}{40-35}$  & \blank{2.5cm} \\
    $[40, 50]$  & $\dfrac{h(50)-h(40)}{50-40}$  & \blank{2.5cm} \\
    $[50, 60]$  & $\dfrac{h(60)-h(50)}{60-50}$  & \blank{2.5cm} \\
    \end{tabularx}

    \item At which interval(s) does the AROC transition from \textbf{positive to negative}?
    What does this indicate about the function's behavior?

    \vspace{4pt}
    \blank{11cm}

    \item At which interval(s) does the AROC transition from \textbf{negative to positive}?

    \vspace{4pt}
    \blank{11cm}

    \item On which portion of the ride is the AROC \textbf{increasing}? What does that tell you
    about the concavity of $h$?

    \vspace{4pt}
    \blank{11cm}

    \item Based on the AROC values, identify a possible \textbf{point of inflection} and explain
    your reasoning.

    \vspace{4pt}
    \blank{11cm}
\end{enumerate}
\end{tcolorbox}

\vspace{0.10in}

\begin{tcolorbox}[colback=white, colframe=black!40,
    title=\textbf{Tier E --- Extend: Degree and Global Extrema},
    fonttitle=\bfseries, arc=2mm, left=3mm, right=3mm, top=2mm, bottom=2mm]

\begin{enumerate}[leftmargin=1.5em, itemsep=8pt]
    \item From the table, count the number of local extrema observed. Based on EK 1.4.A.3
    (between every two distinct real zeros there is at least one local extremum), what is the
    \textbf{minimum possible degree} of $h(t)$? Explain.

    \vspace{4pt}
    \blank{11cm}

    \vspace{2pt}
    \blank{11cm}

    \item The ride begins and ends at $h = 0$, so $t = 0$ and $t = 60$ are zeros. The table
    suggests at least one more zero exists. Assuming the polynomial has degree $n$, does
    $h(t)$ have a global minimum on the domain $[0, 60]$? Use EK 1.4.A.4 in your explanation.

    \vspace{4pt}
    \blank{11cm}

    \vspace{2pt}
    \blank{11cm}
\end{enumerate}
\end{tcolorbox}

\end{document}
